\section{Conclusion}

we conclude...

- generic results vs model dependent

-any observationally ruled out models/important implications?



fdjahdksf

\[\]


\section{Discussion}

\paragraph{constraints}
Observations constrain $n_s \approx 0.97$ with $\sigma \lesssim 0.003$ and an upper bound on the tensor-to-scalar ratio $r \lesssim 0.03$ (95\% C.L.) when combining Planck+BICEP/Keck+ACT+BAO data.\footnote{https://arxiv.org/abs/2512.10613v1}\footnote{https://arxiv.org/html/2507.12459v1}


Extended models (eg, extended Starobinsky\footnote{https://arxiv.org/abs/2312.01010} or alpha-Starobinsky\footnote{https://arxiv.org/abs/2409.05615}) are consistent with current CMB and large-scale structure data, and including reheating uncertainties can widen the allowed parameter space.




\paragraph{Dark Matter Neutrino Interactions}

A recent Nature Astronomy article by Zu et al.\footnote{https://doi.org/10.1038/s41550-025-02733-1} finds a tentative reconciliation to the $S_8$ tension via a nonminimal neutrino-dark matter interaction mechanism. While impressive in its own right, in Supplementary Fig. 4 they analyze the impact of nonminimal $\nu$DM interactions on the $\Lambda$CDM parameters, including the scalar spectral index. 

While the main focus of this article is of course its impact on late time structure homogeneity, there is also an interesting impact on many of the other $\Lambda$CDM parameters: they become much more weakly constrained. While furhter analyses are obviously required to study the more subtly consequences of $\nu$DM interactions, this also opens up more flexibility in beyond-$\Lambda$CDM physics and allows us to reconsider previously discarded models. This is rather unfortunate for us down in the swampland, but is also an important insight into the miracles of possible phyiscs that awaits us.


\todo{write this more coherent:} i think this means that models on the edge of our confidence intervals still have a bit of life? viability of models is famously sensitive to assumptions on late-time cosmology and also microphysics ? so despite the fact $n_s$ is god to inflation, CMB constraints are still model dependent and can be weakened by extensions to the SM and $\Lambda$CDM (eg thru this $\nu$DM extension)

\paragraph{Multifield Extensions}


In the single-field case, the Einstein-frame kinetic term can always be made canonical by a field redefinition $\phi\to\chi$. In multifield NMC inflation, the Einstein-frame kinetic sector instead takes the form of a nonlinear sigma model with an induced field-space metric $G_{IJ}(\phi^K)$ that is generically curved, so one cannot globally “flatten” the kinetic term by simple redefinitions. 

A key qualitative result emphasized by Kaiser is that, despite this extra structure, **conformal stretching of the Einstein-frame potential typically drives an effective single-field attractor**: background trajectories rapidly align, the turn rate tends to remain small, and isocurvature/non-Gaussian signatures are often suppressed in this class of models. The main caveat is EFT control: as Hertzberg stresses, in genuinely multifield settings strong coupling at scales $\sim M_{\rm pl}/\xi$ can appear already in the kinetic/gravity sector, which is why Higgs-inflation interpretations (with multiple scalar degrees of freedom) require extra care.
\[\]



We can very similarly begin with the single field Jordan-frame action except we give the fields indices!

\[
S_J=\int d^4 x \sqrt{-g}\left[\frac{M_P^2}{2} F\left(\phi\right) \R-\frac{1}{2} g^{\mu \nu} \partial_\mu \phi \partial_\nu \phi-V_J\left(\phi\right)\right],
\]
\todo{add indices}

While multifield inflation is in general much more difficult to contend with, we can extract some useful conclusions without diving too far into the math.

\todo{cite sonia paper and other mfi papers}

\todo{fix;robot}
When you introduce multifield extensions, add one sentence explaining what is genuinely new: in multifield NMC, the induced field-space metric $G_{I J}$ is generally curved and cannot be made canonical globally, so the dynamics become a nonlinear sigma model (and this is where isocurvature/turn-rate questions live). That is precisely Kaiser's point around the induced $G_{I J}$ and its consequences


